\section{Threats to Validity}
\label{sec:threats}

\subsection{Construct Validity}

Both datasets encode generic benign/malware labels. Neither establishes
ransomware membership, ransomware behavior, or zero-day status. AVClass metadata
is incomplete and not used as behavioral ground truth. The paper therefore uses
the term Windows malware and treats ransomware-specific detection as future
work.

MalBehavD-V1 prefixes are ordered by API-event count. Event 20 is not 20 seconds,
and a lower arithmetic fusion cost does not imply earlier wall-clock detection.
Detector disagreement is a surrogate fusion signal, not the tracking error of a
physical plant. Terms such as proportional, integral, and derivative describe
the mathematical update only; they do not confer closed-loop stability.

Directional entropy is entropy of positive/negative attribution mass, not
predictive class entropy. Attribution variance is dispersion across coordinates
of one SHAP vector, not epistemic variance. “Explanation-guided” describes the
source of the acquisition signals, not proof that they encode human-understandable
malware mechanisms.

Deletion measures fitted-model behavior under feature replacement. It does not
establish causal features, valid PE transformations, adversarial robustness, or
analyst usefulness. EC, CF, and additivity measure different constructs and are
not aggregated into an Overall Fidelity Score.

\subsection{Internal Validity}

Fusion methods use paired static/dynamic rows and validation-selected settings,
which reduces channel and threshold confounding. Yet the adaptive fusion design
was developed after inspecting the original confirmatory result on the same
1,922-sample cohort. The resulting test-row estimates are development evidence;
same-cohort bootstrap intervals cannot remove method-selection bias. Source
differences between benign and malware data may also influence both channel
scores despite exclusion of explicit verdict and threat-name metadata.

For AL, outer test partitions do not enter query selection or target definition,
and policies share initial sets within a fold/seed. The fixed acquisition weights
and min--max normalization were not selected through an independent outer
sensitivity analysis. Pointwise top-100 selection has no diversity constraint.
The simulated oracle returns retained labels without error, whereas analysts can
abstain, disagree, or use different evidence.

For fidelity, mean and median baselines use training data, avoiding direct test
leakage. Top and random masks are matched by coordinate count but not by
standardized displacement, tree usage, semantic group, or manifold distance.
These differences can inflate the top-minus-random gap. Masked test vectors are
not verified against the PE format.

\subsection{Statistical Conclusion Validity}

Fusion confidence intervals are paired over retained predictions, but only one
development cohort is available. Multiple metrics and method comparisons create
selection opportunities, and no external confirmatory family controls them. FPR
and TPR use validation-selected operating points; the PID--stacking FPR
difference is not a comparison at equal TPR.

The AL comparison unit is the outer fold after averaging five seed trajectories.
Five folds provide effect sizes and direction counts but limited asymptotic
inference. Twenty-five fold--seed runs do not represent 25 independent datasets.
The source report supplies aggregate values and selected paired effects rather
than fold vectors in this checkout, so reconstructed figures omit uncertainty
bars.

Restricted-FPR estimates use small test folds. With approximately 300 benign
rows, one false positive changes FPR by roughly 0.33 percentage point. Results at
1\% or 5\% caps cannot validate alert rates at $10^{-3}$ or $10^{-4}$.

The 25 fidelity models share folds and data; they are not independent population
draws. Random masks characterize conditional mask variability. Pairwise
stability values share fitted models and cannot be entered as independent rows
in a naive significance test.

\subsection{External Validity}

The 2,000-row EMBER subset covers a short 2018 period and uses random row folds.
It does not impose temporal, family, campaign, or source separation. The
MalBehavD-V1 cohort is also one source with metadata-availability and minimum-
event selection. Results can change for newer software, different organizations,
other sandbox environments, or another feature extractor.

The horizon-40 subset was selected after horizon-20 results were known. Robust
PID v2 and controlled stress findings are therefore explicitly exploratory. The
timestamp pilot contains only eight usable sequences and cannot support
performance inference. No measured result covers production queueing, analyst
workflow, or end-to-end incident response.

LightGBM and TreeSHAP are central to RQ2/RQ3. The acquisition summaries can be
defined for other additive attributions, but their numeric behavior, cost, and
stability may differ for raw-byte networks, graph models, or sequence models.

\subsection{Integration Validity}

The three RQs use two datasets and different feature spaces. The architecture
connects them as stages of an auditable learning process, but the experiments do
not demonstrate a single deployed pipeline. There is no measured gate deciding
which EMBER sample receives dynamic analysis, no AL retraining of the
MalBehavD-V1 fusion channels, and no shared latency/cost objective. Claims about
synergy, interaction, or complete-system superiority are outside the evidence.

\subsection{Reproducibility Validity}

The source-of-truth report names commit \texttt{8bfa06b65b93}, dataset hashes,
result directories, and verification artifacts. The remote branch available in
the writing environment ends at an older commit, so the manuscript was generated
from the supplied report rather than by rerunning the missing artifacts. Tables
and plots reproduce reported aggregate values; they do not constitute an
independent computational replication.

Before submission, the named commit and all result directories must be pushed,
hashes checked, verification commands rerun in a clean environment, and tables
and figures regenerated directly from fold-level records. Bibliographic metadata,
data-sharing rights, authorship order, corresponding author, and the target
journal template also require final verification. Computational replay can show
consistency with stored code and data; it cannot prove upstream label accuracy or
external generalization.
